\section{Experimental setup}
\label{sec:setup}

\subsection{Training data}
\label{sec:data}

We train on the Reason+Long mixture of LLM-as-Jev \citep{li2026llmasjev}, with Nimble added: 13 datasets covering intent routing (CLINC150 \citep{clinc150}, MASSIVE \citep{massive}, Bitext \citep{bitext}), short reasoning (ReClor \citep{reclor}, LogiQA~2.0 \citep{logiqa2}, CosmosQA \citep{cosmosqa}, PIQA \citep{piqa}, ANLI \citep{anli}, ProofWriter \citep{proofwriter}), long input (ContractNLI \citep{contractnli}, MuSiQue \citep{musique}, QuALITY \citep{quality}) and Nimble \citep{nimble}.
Each row is a state (the passages), a question and a list of options.
For example, a MuSiQue row gives about twenty paragraphs and a multi-hop question, and asks whether the paragraphs contain enough information to answer it (``Yes'' or ``No'').

We keep \NPerDataset{} training rows per dataset, \NTrain{} in total, after these filters:
\begin{enumerate}\setlength{\itemsep}{0pt}\setlength{\topsep}{2pt}
  \item \textbf{Duplicates.} Rows with the same input are collapsed, and removed if their answers differ.
  \item \textbf{Leakage.} A training row is removed if its input or source document appears in the validation split.
  \item \textbf{Benchmark overlap.} A training row is removed if it shares a 13-word n-gram, or a sentence of six or more words, with an evaluation benchmark.
  \item \textbf{Length.} Prompts over \MaxTokens{} tokens (Qwen3.5-4B tokenizer) are dropped, never truncated.
\end{enumerate}

The held-out rows come from the validation split, with the same duplicate and length filters, up to 500 per dataset and in whole source documents (\NHeldOut{} rows).
A random 100 rows per dataset form the dev set (\NDev{} rows), used to choose learning rates and the evidence threshold; the other \NTest{} rows form the test set, excluded from model selection.
The split is by row, so a dev and a test row can share a source document, such as two statements about one contract.

\subsection{Models}
\label{sec:models}

UltraBERT-380M \citep{ultrabert} is fully fine-tuned with each of the four recipes of \cref{sec:loss}.
Qwen3-0.6B \citep{qwen3} is fully fine-tuned with the decoder loss, with a frozen copy as $p_0$; it is the matched comparison.
As references, we also report two LoRA models \citep{lora}, Qwen3-0.6B and Qwen3-4B, reproduced from LLM-as-Jev on the same training set and loss.

\subsection{Implementation details}
\label{sec:Implementation details}

Every model is trained for one epoch with a batch of 32 rows.
We sweep the learning rate on dev, one run per value: $\{0.8, 1.5, 3, 5\}\times10^{-5}$ for the encoder with the softmax over ``yes'', $\{0.8, 1.5, 2, 3\}\times10^{-5}$ with the sigmoid-normalised loss, and $\{0.5, 1, 2\}\times10^{-5}$ for the decoder.
Both encoder sweeps choose $1.5\times10^{-5}$, which we reuse for the other two recipes; for Qwen the two best rates differ by 0.1 points, and we use $1\times10^{-5}$.
At the chosen rate we train each model several times, changing only the row order and option shuffling, and report the average.

The metric is accuracy, and ALL is its mean over the 13 datasets.
On the test set we average each row's correctness over a model's runs before computing ALL.
For a difference between two models we give a paired 95\% bootstrap interval: 2{,}000 resamples of the test rows within each dataset, with both models scored on the same rows.
It covers test-row sampling only, not retraining, and treats rows from one document as independent, so it can be too narrow.

The encoder is scored with fp32 weights under bf16 autocast, since bf16 weights change its long-document results (\cref{sec:precision}); the Qwen checkpoints use bf16 weights.

\subsection{External benchmarks and cost}
\label{sec:external}

To see how training changes behaviour outside the mixture, we score six benchmarks as decision rows, never used for any choice: BoolQ \citep{boolq}, MMLU \citep{mmlu}, MMLU-Pro \citep{mmlupro}, ARC-Challenge \citep{arc}, Winogrande \citep{winogrande} and SciQ \citep{sciq}.
We draw fixed random samples of 500, 1200, 800, 600, 500 and 400 rows, the sizes of LLM-as-Jev, whose samples are not released.
External is the mean accuracy over the six.

Time per decision is the mean latency of one request over the dev rows, scored one at a time on one A100, as when a pipeline sends each item as it arrives.
It covers the forward pass and option scoring but not tokenisation, after 20 warm-up rows, with the GPU synchronised around every call.
The dev rows average 1{,}100 to 1{,}200 tokens.
The encoder is timed on the softmax checkpoint and the decoders on the reproduced LoRA checkpoints with merged weights; the other checkpoints of each model share their architecture.
\Cref{app:impl} gives batched throughput.
